\section{Executive Summary}
\label{sec:exec-summary}

This program proves two customer adoption architectures by November~30, 2026, on one common 10BASE-T1S electrical infrastructure. They are not competing product strategies -- they are two ways a customer enters the same Alfred network.

\begin{diagram}
                        ALFRED

                    10BASE-T1S
                         |
             +-----------+-----------+
             |                       |
             v                       v

      CLEAN-SHEET                 BROWNFIELD
      NEW PROGRAM              EXISTING PROGRAM
             |                       |
       direct E2B                 Gateway
       endpoints                    |
             |                 CAN/CAN-FD/LIN
             v                       |
       peripherals               existing ECUs
\end{diagram}

\textbf{Architecture~A (Greenfield):} a clean-sheet customer adopts Alfred's Ethernet-native infrastructure directly. Peripherals connect through \textbf{LAN8660} (Control Endpoint) and \textbf{LAN8661} (Lighting Endpoint) -- Microchip's MCU-less, RCP-managed 10BASE-T1S endpoint devices -- with no Alfred-authored firmware running at the edge node at all.

\textbf{Architecture~B (Brownfield):} an existing-vehicle customer keeps their CAN/CAN-FD/LIN network exactly as it is. An Alfred Gateway (a customer MCU running Alfred's own capture/timestamp/mirror/translate logic, placed onto 10BASE-T1S via \textbf{LAN8651}) attaches to that network and represents it on an Alfred-facing pseudo-bus.

\begin{decision}[title={One platform philosophy, three sibling node classes -- not one interchangeable board}]
LAN8660 and LAN8661 are \textbf{mandatory} for the November Greenfield proof -- they are Alfred's actual Control and Lighting E2B endpoints, not a stretch goal. They are also MCU-less, fixed-function devices, while the Gateway role structurally requires a programmable MCU running Alfred's own capture/translate logic. Because of that, \textbf{one identical PCB cannot become both an E2B and a Gateway} by component population -- the two roles need genuinely different silicon. The architecture instead shares one \emph{electrical platform philosophy} (Section~\ref{sec:platform-strategy}) across three sibling node classes -- E2B-Control (LAN8660), E2B-Lighting (LAN8661), and Gateway (MCU~+~LAN8651) -- common power, protection, connectors, mechanical language, and validation discipline, not common silicon.
\end{decision}

\textbf{LAN8651} is the host-side MAC-PHY used wherever a programmable MCU needs 10BASE-T1S access: the Clicker~4 (temporary central-compute stand-in, out of scope for redesign) and the Brownfield Gateway MCU both reach the network through it. Commercial LAN8651 hardware (MikroE Two-Wire ETH Click) already does this job well -- the program does not spend its first PCB revision reproducing it.

The November~30 exit gate (Section~\ref{sec:exit-criteria}) requires, on real hardware, with measured data: a Clicker~4 and at least one LAN8660 and one LAN8661 coexisting on a multidrop PLCA segment, LAN8660 driving a real motor-driver/actuator path with observable telemetry, LAN8661 driving a real lighting load, both simultaneously, and, separately, a Gateway passively mirroring a live CAN network onto the same T1S infrastructure and then, in an explicitly authorized mode, carrying a command from the Clicker~4 through to a physical device and back -- all against the quantitative validation discipline in Section~\ref{sec:measurement}, never against a qualitative ``it worked.''
